
\subsubsection{Findings}

The primary finding is that Marchenko-Pastur-defined curvature subspace alignment differentiates ERM from robust training with large effect size (Δ = 40.8 percentage points, Cohen's d = 1.87) in a single-seed experiment on Waterbirds. This geometric signature is accompanied by curvature concentration (53\% more outlier eigenvalues in ERM) and measurable SGD directional bias (0.15, p = 0.023).

Two mechanism components are incomplete. Minority gradient alignment to real Hessian eigenvectors was not tested due to use of a random basis, leaving a gap between curvature concentration and SGD directional bias. Geometry-phenotype coupling was not observed along interpolation paths (ρ = 0.045), suggesting that geometric orientation may classify training regimes without predicting robustness within a regime, or that the 5-epoch training was insufficient to reveal coupling.

\subsubsection{Limitations}

The results are limited to Waterbirds with one random seed for the primary experiment. Generalization to other datasets (CelebA, Colored MNIST) and other spurious correlation types (attribute-based rather than visual background) has not been tested. The single-seed result demonstrates feasibility with large effect size but does not establish statistical robustness across multiple training runs.

We used a random orthonormal basis rather than computing actual Hessian eigenvectors for the minority alignment test, leaving this mechanism component untested rather than validated or refuted. Computing the top-23 eigenvectors from h-m1 validation would require 4-6 hours but was not performed.

Experiments were limited to ResNet-50. Architecture invariance across ViT, Wide ResNet, or shallow CNNs has not been tested. The Marchenko-Pastur method is designed for over-parameterized networks and may not apply to under-parameterized regimes.

Interpolation path experiments used only 5-epoch training, producing poor endpoint worst-group accuracy differentiation (ERM: 32.3\%, DRO: 9.8\%). Full training to convergence (100 epochs) was not performed for these experiments, limiting conclusions about geometry-phenotype coupling.

The diagnostic framework requires minority group identification for gradient computation, though not for model training. This distinguishes it from fully unsupervised detection methods while avoiding annotation costs during training.

\subsubsection{Interpretation}

The failure to observe geometry-phenotype coupling along paths (ρ = 0.045) has three potential explanations. First, coupling may emerge only at converged solutions, not along training trajectories or interpolation paths. The 5-epoch training was insufficient to produce converged endpoints. Second, geometric alignment may reliably classify training regimes (ERM vs. DRO) without predicting worst-group accuracy within a regime, functioning as a binary classifier rather than a continuous robustness predictor. Third, linear and Fast Geometric Ensembling interpolation may not preserve functional relationships present in actual training trajectories.

The validated components (geometric signature, curvature concentration, SGD bias) function independently of the coupling result. The geometric signature exists and is explained by curvature concentration plus SGD bias, regardless of whether it predicts worst-group accuracy along arbitrary paths.
